\clearpage\thispagestyle{empty}
\pdfbookmark[0]{Sources for Geocentric Positions of the Planets (1804)}{bemerkungen1804-sources}
\begin{center}\textbf{Sources for Geocentric Positions of the Planets (1804)}\end{center}
\noindent This AI translation follows Carl Friedrich Gauss’s German text in \textit{Werke}, vol.~VI (Göttingen, 1874), pp.~94--105, including the handwritten remarks printed at the end. The paper first appeared as article L in the May 1804 issue of \textit{Monatliche Correspondenz zur Beförderung der Erd- und Himmels-Kunde}, edited by Franz Xaver von Zach, vol.~9, pp.~385--400. \href{https://gdz.sub.uni-goettingen.de/id/PPN236007467}{Collected German text}. \href{https://www.digitale-sammlungen.de/en/view/bsb10538601?page=421}{First printing}.
\par\smallskip\noindent The collected text gives \(x-X\) as the denominator in the formula for \(\operatorname{tang}\alpha\) (p.~96). The first printing has \(x-Z\) (p.~387), as does the reprint in Appendix VII, pp.~53--58, of \href{https://gdz.sub.uni-goettingen.de/id/PPN60187157X}{Carl Haase’s German translation of \textit{Theoria motus}} (Hannover, 1865). In the collected handwritten remarks, \textit{in plano ipsius} \(Z\) means “in the plane of \(Z\).”
\par\smallskip\noindent Gauss treats the related problem in Articles~53--57, pp.~60--66, of \href{https://www2.karlin.mff.cuni.cz/~tuma/Aplikace21/theoryofmotionof00gausuoft.pdf#page=89}{\textit{Theoria motus}, translated by Charles Henry Davis} (Boston, 1857). Article~56, p.~65, refers back to the May 1804 paper as another solution.
\par\smallskip\noindent Translation and typesetting produced by AI.
\clearpage
